\section{Обзор существующих решений}
\label{sec:Chapter2} \index{Chapter2}

В разделе рассматриваются подходы к улучшению качества BVH. Для этого систематизация ведётся по двум осям: (i)~уровень
вмешательства в pipeline (внутрь билдера, на входе билдера, поверх
готового BVH) и (ii)~применимость в условиях, диктуемых
DXR/Vulkan~RT API. Для каждого подхода критически оценивается,
насколько он удовлетворяет ограничениям, сформулированным в
постановке задачи: неизменность кода
вендорского билдера, сохранение визуальной эквивалентности,
аппаратные лимиты на 16-битные индексы и число
geometry-описаний на BLAS.

\subsection{Классические top-down SAH-билдеры}

Поверхностно-площадная эвристика SAH
была введена Goldsmith и Salmon в 1987~г.~\cite{goldsmith1987} и
формализована для контекста ray-tracing-обхода
MacDonald и Booth в 1990~г.~\cite{macdonald1990}. SAH постулирует,
что вероятность попадания случайного луча в узел $X$ оценивается
отношением площадей поверхности,
$P(\text{hit}\,X)\!\approx\!SA(X)/SA(\text{root})$,
откуда выводится ожидаемая стоимость обхода всего дерева.

Top-down SAH-билдер на каждом узле перебирает кандидатов разбиения и выбирает $\arg\min$
SAH-стоимости. Точный перебор дорог на больших наборах примитивов.
Wald~\cite{wald2007} показал, что достаточно вдоль каждой оси
пространство делить на $16$~равных интервалов, примитивы
распределяются по ним, а лучшая плоскость разбиения выбирается среди
границ этих интервалов. Такой вариант работает за $O(n\log n)$ и даёт качество BVH,
отличающееся от варианта полного перебора менее чем на $1\,\%$, поэтому стал
стандартом в индустрии.

\textit{Соответствие задаче.} Этот подход --- \emph{внутренний} алгоритм
билдера. В DXR/Vulkan~RT он работает внутри драйвера, и для
разработчика доступа к нему нет. При этом любая SAH-оптимизация внутри
билдера улучшает результат лишь в рамках заданного входа: если
поданный на вход материальный подмеш имеет AABB размером со всю сцену, SAH-разбиение внутри него не
может предотвратить появление огромного BLAS, который в дальнейшем приведёт к перекрытиям в TLAS.

\subsection{HLBVH и LBVH: параллельная сборка по Morton-кривой}

Lauterbach и др. предложили LBVH (Linear~BVH): кодирование
центроидов примитивов 30-битным Morton-кодом, радикс-сортировка по
этому коду и рекурсивное построение дерева по битам ключа.
Pantaleoni и Luebke расширили подход до HLBVH~\cite{HLBVH},
реализовав параллельную постройку BVH для динамических ассетов
на GPU.

Идея Morton-сортировки используется и независимо от LBVH:
предварительная перестановка примитивов по Z-кривой
улучшает локальность
доступа в кэше и обычно даёт билдеру более стабильное разбиение. Это поведение закреплено как baseline-приём,
упоминаемый в практических руководствах
NVIDIA~\cite{nvOptMesh} и AMD~\cite{rraManual}.

\textit{Соответствие задаче.} HLBVH как замена вендорского билдера
не применим: для DXR/Vulkan-сцен дерево обязан построить драйвер.
Однако \emph{глобальная Morton-перестановка примитивов} на стадии,
предшествующей вызову билдера, полностью совместима с
DXR-контрактом, и именно этот приём используется в
настоящей работе как один из двух обязательных шагов
предлагаемого алгоритма: после
объединения треугольников всех материальных групп они
сортируются по Morton-коду центроида перед разбиением на
кластеры.

\subsection{SBVH и SAH-guided pre-splitting}

Параллельно линии классических билдеров развивалась идея ослабления
требования «один треугольник --- один лист». Stich и
др.~\cite{stich2009} (метод SBVH) показали, что разрешение
разрезать треугольник на несколько ссылок
повышает качество дерева на 10--30\,\% по SAH-стоимости при
умеренном раздувании числа ссылок. Ganestam и
Doggett~\cite{ganestam16_splitting} перенесли эту идею на стадию
\emph{пре-разрезания} больших треугольников до запуска top-down
билдера, что ускоряет последующую постройку и сохраняет
SBVH-качество.

\textit{Соответствие задаче.} Обе работы атакуют ту же
фундаментальную проблему --- большой AABB примитива, портящий
верхние узлы дерева, --- но на уровне отдельных треугольников. SBVH
встраивается в билдер и поэтому в DXR-сцене недоступен;
pre-splitting Ganestam--Doggett в принципе реализуем на стороне
CPU, однако он кратно увеличивает число треугольников меша, что в
условиях многократного использования ассета (рендер-проход, тени,
GI и т.\,п.) даёт значительную нагрузку на остальной pipeline.
Предлагаемый в настоящей работе подход концептуально близок к
pre-splitting на уровне идеи «компактизация на входе билдера», но
оперирует не отдельными треугольниками, а \emph{целыми крупными ассетами}:
исходный меш не разбивается на части внутри одного BLAS, а перегруппировывается
по пространственной близости в независимые мини-ассеты, и тем самым решает проблему перекрытий в TLAS.

\subsection{Метрики качества BVH}

Стандартная SAH-стоимость удобна, но является лишь \emph{прокси}
для реального времени трассировки. Aila, Karras и
Laine~\cite{ailaKarrasLaine2013} провели систематический анализ
корреляции офлайн-метрик с GPU-измерениями и показали:
(i)~на 20+ тестовых сценах коэффициент детерминации SAH cost
против реального времени \texttt{DispatchRays} составил
$R^2\!\approx\!0.85$;
(ii)~альтернативная метрика EPO (Expected~Path~Overlap) улучшает
корреляцию, но вычислительно дороже ($O(n^2)$ в худшем случае);
(iii)~кэш-эффекты и SIMD-ширина дают расхождение SAH-время до
15\,\% на сценах с регулярной структурой.

Karras и Aila~\cite{karras2013} также показали, что для
GPU-friendly~BVH важна не только SAH-стоимость, но и
сбалансированность дерева; несбалансированные SBVH-деревья могут
уступать сбалансированным LBVH-деревьям при одинаковой
SAH-стоимости.

\textit{Соответствие задаче.} Поскольку в DXR/Vulkan-сцене дерево
строится драйвером и реальная стоимость может зависеть от
драйверо-специфичных эвристик, в настоящей работе в качестве
основной метрики выбрано \emph{измеренное GPU-время} основного
прохода трассировки теней \texttt{rtsm::do\_trace}. Эта метрика
непосредственно отражает практический эффект оптимизации и не
зависит от выбранной офлайн-метрики; SAH cost упоминается лишь
как теоретическое объяснение причины ускорения.

\subsection{Альтернативные ускоряющие структуры}

Помимо BVH рассматриваются и альтернативные подходы. Aman и
др.~\cite{tetra2021} предложили компактное представление сцены на
основе тетраэдрализации, дающее меньший объём данных и линейную
сложность поиска ближайшего пересечения для протяжённых лучей.
KD-деревья и uniform~grid традиционно используются в
офлайн-рендерах. NVIDIA отдельно использует \emph{rebraiding}
TLAS-инстансов для динамической перестройки верхних
уровней~\cite{nvOptMesh}.

\textit{Соответствие задаче.} Все альтернативы предполагают замену
ускоряющей структуры целиком, что невозможно в условиях
фиксированного DXR-pipeline. Они рассматриваются как ориентиры для
понимания нижней границы возможной стоимости трассировки, но не
как конкуренты предлагаемого метода.

\subsection{Промышленные руководства и анализаторы}

NVIDIA в техническом руководстве~\cite{nvOptMesh} прямо рекомендует
разработчикам ассетов: (i)~избегать ассетов с диагональю
AABB, сопоставимой с размером всей сцены; (ii)~разрезать «жирные»
ассеты вручную в DCC-инструментах; (iii)~контролировать число
geometry-описаний на BLAS, желательно не более 16--32. AMD в
руководстве к RRA~\cite{rraManual} описывает аналогичный набор
рекомендаций и предоставляет инструмент для визуализации
SAH-стоимости и overlap-а во внутренних узлах BLAS непосредственно
из захвата GPU-памяти.

Спецификация DXR~\cite{dxrSpec} фиксирует контракт интерфейса:
поля \texttt{geometry\_index} и \texttt{primitive\_index} должны
оставаться стабильными при обходе и в hit-шейдере; диапазон
индексов одной \texttt{GeometryDesc} является \emph{непрерывным}. Это формальное
обоснование того, почему драйвер не может «разрезать» большую
geometry самостоятельно.

\textit{Соответствие задаче.} Эти руководства подтверждают практическую
актуальность проблемы, описанной в постановке. Однако они
оставляют выполнение рекомендаций за разработчиком ассетов и не
предлагают автоматизированного инструмента, работающего на
стороне runtime игрового движка. Настоящая работа закрывает эту
нишу: предлагается алгоритм и его реализация, которые выполняют
рекомендации NVIDIA автоматически при инициализации BVH-контекста
без вмешательства в ассет-пайплайн.

\subsection{Сводное сравнение}

В сводной таблице сведены ключевые свойства
рассмотренных подходов с позиции ограничений настоящей работы.

\begin{table}[H]
\centering
\caption{Сравнение подходов по совместимости с DXR-pipeline}
\label{tab:related_work}
\small
\begin{tabular}{lccc}
\hline
\textbf{Подход} & \textbf{Уровень} & \textbf{Совместим} & \textbf{Решает перекрытия} \\
& \textbf{вмешательства} & \textbf{с DXR} & \textbf{в TLAS} \\
\hline
Top-down SAH binner~\cite{wald2007} & внутри билдера & --- & частично \\
HLBVH~\cite{HLBVH} & замена билдера & нет & нет \\
Morton-reorder & вход билдера & да & нет \\
SBVH~\cite{stich2009} & внутри билдера & --- & да \\
SAH pre-split~\cite{ganestam16_splitting} & вход билдера & да$^*$ & да \\
Tetra-AS~\cite{tetra2021} & замена структуры & нет & --- \\
\textbf{Глобальная пространственная} & \textbf{вход билдера} & \textbf{да} & \textbf{да} \\
\textbf{кластеризация (наш)} & & & \\
\hline
\end{tabular}

\vspace{0.5em}
\raggedright\footnotesize
$^*$~требует дисциплины: pre-split может вылезти за лимит числа
геометрий или 16-битного индекса, если применяется к большим
группам без учёта аппаратных ограничений.
\end{table}

\subsection{Выводы по обзору}

\begin{enumerate}
    \item Существующие алгоритмы построения BVH (top-down
      SAH binner, HLBVH, SBVH) показывают, что классические
      эвристики хорошо справляются с улучшением дерева,
      \emph{если} на вход подан пространственно компактный
      набор примитивов.
    \item Реальный DXR/Vulkan-pipeline ставит билдер в условия
      чёрного ящика и фиксирует контракт непрерывности
      индекс-диапазонов и формат 16-битных индексов. Любое
      решение, требующее модификации внутри билдера, этим
      ограничениям не удовлетворяет.
    \item Доступная разработчику зона вмешательства --- стадия до
      подачи в билдер. На этой стадии существуют отдельные приёмы
      (Morton-reorder, pre-splitting), но ни один из них не
      закрывает сценарий «отдельные материальные подмеши имеют размер всей сцены, вызывая огромные перекрытия в TLAS» при
      сохранении контракта DXR.
    \item Все упомянутые работы оценивают результат либо
      SAH-стоимостью, либо время-симулятором, но не на реальном
      production-движке. В настоящей работе оценка ведётся
      непосредственно по GPU-времени трассировки теней в
      работающем семпле DagorEngine, что повышает достоверность
      результата.
\end{enumerate}

Эти выводы определяют как цель алгоритмической части работы
при описании алгоритма, так и форму экспериментальной
валидации.
